\documentclass{tikzposter}
\usepackage[utf8]{inputenc}
\usepackage[T1]{fontenc}
\usepackage[french]{babel}
\usepackage{amsmath,amssymb}
\usepackage{graphicx}
\usepackage{hyperref}

\title{Makeitalive: Animation de Paysages Statiques via Generative AI}
\author{Arthur et Equipe}
\institute{Projet de Recherche en IA Générative}

\begin{document}
\maketitle

\begin{columns}
\begin{column}{0.33}
\block{Contexte du Projet}{
Le projet \textbf{Makeitalive} vise à animer des images statiques de paysages en utilisant des techniques d'Intelligence Artificielle Générative. L'objectif principal est de donner vie à des photographies de paysages en générant un mouvement naturel et fluide.

\textbf{Problématique:} Comment transformer une image fixe en une séquence animée réaliste ?
}

\block{Collection de Données}{
Nous avons collecté des données en extrayant des vidéos YouTube de paysages :
\begin{itemize}
\item Téléchargement automatique de vidéos haute qualité
\item Extraction d'images consécutives (frames)
\item Création de paires d'images pour l'entraînement
\item Dataset final : milliers de paires img\_A $\rightarrow$ img\_B
\end{itemize}

\textbf{Commandes utilisées:}
\begin{verbatim}
uv run src/data/download_youtube.py \
    --url "https://..." \
    --out "./data/10hours.mp4"

uv run src/data/make_dataset_video.py \
    --video "./data/10hours.mp4" \
    --name "dataset_local"
\end{verbatim}
}
\end{column}

\begin{column}{0.33}
\block{Approches Techniques}{
Nous avons implémenté deux stratégies principales :

\textbf{1. Fine-tuning LoRa de Modèles Pré-entraînés}
\begin{itemize}
\item Utilisation de modèles de diffusion stables
\item Adaptation légère via LoRa pour réduire la complexité
\item Focus sur la génération conditionnée
\end{itemize}

\textbf{2. Entraînement from Scratch de Motion Flow}
\begin{itemize}
\item Architecture U-Net personnalisée
\item Prédiction directe du champ de mouvement
\item Entraînement end-to-end sur notre dataset
\end{itemize}

\textbf{Architecture MotionFlowUNet:}
\begin{itemize}
\item Encodeur: 4 niveaux de downsampling
\item Bottleneck: 1024 canaux
\item Décodeur: Upsampling avec skip connections
\item Sortie: Champ de mouvement 2D (dx, dy)
\end{itemize}
}

\block{Résultats Préliminaires}{
\begin{itemize}
\item Modèle entraîné sur 100 époques
\item Loss MSE sur le flow prédit
\item Génération de vidéos à partir d'images statiques
\item Visualisation des champs de mouvement
\end{itemize}

\textbf{Métriques d'évaluation:}
\begin{itemize}
\item Précision du flow prédit
\item Qualité visuelle des animations
\item Cohérence temporelle
\end{itemize}
}
\end{column}

\begin{column}{0.33}
\block{Impact et Applications}{
\textbf{Applications potentielles:}
\begin{itemize}
\item Animation automatique de photographies
\item Génération de contenu pour médias
\item Outils créatifs pour artistes
\item Recherche en vision par ordinateur
\end{itemize}

\textbf{Avantages de l'approche:}
\begin{itemize}
\item Pas de besoin de vidéos sources complexes
\item Génération de mouvement naturel
\item Contrôle fin du mouvement via le modèle
\end{itemize}
}

\block{Prochaines Étapes}{
\begin{itemize}
\item Amélioration de la qualité des animations
\item Intégration de conditions (direction, vitesse)
\item Extension à d'autres types d'images
\item Comparaison quantitative des approches
\end{itemize}

\textbf{Défis techniques:}
\begin{itemize}
\item Génération de mouvement réaliste
\item Évitement des artefacts visuels
\item Optimisation des performances
\end{itemize}
}
\end{column}
\end{columns}

\end{document}